% GENERATED FILE. Edit canonical lessons, then run npm run generate:books.
% canonical-source-hash: fnv1a64:1607770e882050b3
% canonical-lessons: MW-C148-buro, MW-C148-jaruri, MW-C148-khas, MW-C148-sado, MW-C148-anokho

\chapter{Wicked, Necessary, Special, Plain, Unusual}
\label{ch:mw-wicked-necessary-special-plain-unusual}
\hlchaptermodality{148}

\begin{chapteropening}
\textbf{By the end of this chapter:} \emph{I can say wicked, necessary, special, plain and unusual.}

Complete the last lesson of chapter 148: I can say wicked, necessary, special, plain and unusual.
\end{chapteropening}

\section[\texorpdfstring{\mw{बुरो} (buro)}{बुरो (buro)}]{\mw{बुरो} (buro) — wicked, bad}

\label{lesson:MW-C148-buro}



\begin{quote}
Before the new one: say the Marwadi for miserly, then the Marwadi for good-natured.
\end{quote}

\subsection*{You'll want to know: \mw{बुरो}}
\textbf{\mw{बुरो}} — \emph{buro} — \textquotedblleft{}wicked, bad\textquotedblright{}.

\begin{cousinweb}
\textbf{\mw{बुरो}} (\emph{buro}) is a tadbhava, a word inherited and reshaped by sound change. It comes from \emph{*bura}, \textquotedblleft{}defective\textquotedblright{}, a form CDIAL reconstructs from its descendants. CDIAL gives no Marwari form in this entry (CDIAL 9289); the match is through Hindi \emph{burā} and Gujarati \emph{būrũ}. CDIAL sets it in a large group of \textquotedblleft{}defective\textquotedblright{} words.
\end{cousinweb}

\begin{grammarlens}[title={What you've built}]
One more word for this chapter.
\end{grammarlens}

\subsection*{Your turn}
\begin{itemize}
\raggedright
  \item \emph{Say it:} \emph{buro}
  \item \emph{Say it:} \emph{buro}, once more
  \item \emph{Say it:} say \emph{buro}
  \item \emph{Recall:} say the Marwadi for miserly, then the Marwadi for good-natured, then say \emph{buro} again
\end{itemize}

\begin{tcolorbox}[breakable,colback=teal!4,colframe=teal!35!black,title={Before you move on}]
What does \mw{बुरो} mean? (\textquotedblleft{}Wicked, bad\textquotedblright{}.) Say it once more.
\end{tcolorbox}
\section[\texorpdfstring{\mw{जरूरी} (jarūrī)}{जरूरी (jarūrī)}]{\mw{जरूरी} (jarūrī) — necessary}

\label{lesson:MW-C148-jaruri}



\begin{quote}
Before the new one: say the Marwadi for good-natured, then the Marwadi for wicked.
\end{quote}

\subsection*{You'll want to know: \mw{जरूरी}}
\textbf{\mw{जरूरी}} — \emph{jarūrī} — \textquotedblleft{}necessary\textquotedblright{}.

\begin{cousinweb}
\textbf{\mw{जरूरी}} (\emph{jarūrī}) is a loan from Arabic \emph{ḍarūrī}, \textquotedblleft{}necessary\textquotedblright{}, which reached Marwadi through Persian. Its Arabic root is \emph{ḍ-r-r} (Lane). Persian says it \emph{zarūrī}; Marwadi says the \emph{z} as \emph{j}.
\end{cousinweb}

\begin{grammarlens}[title={What you've built}]
One more word for this chapter.
\end{grammarlens}

\subsection*{Your turn}
\begin{itemize}
\raggedright
  \item \emph{Say it:} \emph{jarūrī}
  \item \emph{Say it:} \emph{jarūrī}, once more
  \item \emph{Say it:} say \emph{jarūrī}
  \item \emph{Recall:} say the Marwadi for good-natured, then the Marwadi for wicked, then say \emph{jarūrī} again
\end{itemize}

\begin{tcolorbox}[breakable,colback=teal!4,colframe=teal!35!black,title={Before you move on}]
What does \mw{जरूरी} mean? (\textquotedblleft{}Necessary\textquotedblright{}.) Say it once more.
\end{tcolorbox}
\section[\texorpdfstring{\mw{खास} (khās)}{खास (khās)}]{\mw{खास} (khās) — special}

\label{lesson:MW-C148-khas}



\begin{quote}
Before the new one: say the Marwadi for wicked, then the Marwadi for necessary.
\end{quote}

\subsection*{You'll want to know: \mw{खास}}
\textbf{\mw{खास}} — \emph{khās} — \textquotedblleft{}special\textquotedblright{}.

\begin{cousinweb}
\textbf{\mw{खास}} (\emph{khās}) is a loan from Arabic \emph{ḫāṣṣ}, \textquotedblleft{}particular, special\textquotedblright{}, which reached Marwadi through Persian. Its Arabic root is \emph{ḫ-ṣ-ṣ} (Lane).
\end{cousinweb}

\begin{grammarlens}[title={What you've built}]
One more word for this chapter.
\end{grammarlens}

\subsection*{Your turn}
\begin{itemize}
\raggedright
  \item \emph{Say it:} \emph{khās}
  \item \emph{Say it:} \emph{khās}, once more
  \item \emph{Say it:} say \emph{khās}
  \item \emph{Recall:} say the Marwadi for wicked, then the Marwadi for necessary, then say \emph{khās} again
\end{itemize}

\begin{tcolorbox}[breakable,colback=teal!4,colframe=teal!35!black,title={Before you move on}]
What does \mw{खास} mean? (\textquotedblleft{}Special\textquotedblright{}.) Say it once more.
\end{tcolorbox}
\section[\texorpdfstring{\mw{सादो} (sādo)}{सादो (sādo)}]{\mw{सादो} (sādo) — plain, simple}

\label{lesson:MW-C148-sado}



\begin{quote}
Before the new one: say the Marwadi for necessary, then the Marwadi for special.
\end{quote}

\subsection*{You'll want to know: \mw{सादो}}
\textbf{\mw{सादो}} — \emph{sādo} — \textquotedblleft{}plain, simple\textquotedblright{}.

\begin{cousinweb}
\textbf{\mw{सादो}} (\emph{sādo}) is a loan, usually traced to Persian \emph{sāda}, \textquotedblleft{}plain, simple\textquotedblright{}.
\end{cousinweb}

\begin{grammarlens}[title={What you've built}]
One more word for this chapter.
\end{grammarlens}

\subsection*{Your turn}
\begin{itemize}
\raggedright
  \item \emph{Say it:} \emph{sādo}
  \item \emph{Say it:} \emph{sādo}, once more
  \item \emph{Say it:} say \emph{sādo}
  \item \emph{Recall:} say the Marwadi for necessary, then the Marwadi for special, then say \emph{sādo} again
\end{itemize}

\begin{tcolorbox}[breakable,colback=teal!4,colframe=teal!35!black,title={Before you move on}]
What does \mw{सादो} mean? (\textquotedblleft{}Plain, simple\textquotedblright{}.) Say it once more.
\end{tcolorbox}
\section[\texorpdfstring{\mw{अनोखो} (anokho)}{अनोखो (anokho)}]{\mw{अनोखो} (anokho) — unusual}

\label{lesson:MW-C148-anokho}



\begin{quote}
Before the new one: say the Marwadi for special, then the Marwadi for plain.
\end{quote}

\subsection*{You'll want to know: \mw{अनोखो}}
\textbf{\mw{अनोखो}} — \emph{anokho} — \textquotedblleft{}unusual\textquotedblright{}.

\begin{cousinweb}
\textbf{\mw{अनोखो}} (\emph{anokho}) is a tadbhava, a word inherited from Sanskrit and reshaped by sound change. It comes from \emph{anapekṣa}, \textquotedblleft{}careless\textquotedblright{}. CDIAL gives no Marwari form in this entry (CDIAL 289); the match is through Hindi \emph{anokhā}. It is built on \emph{an-}, \textquotedblleft{}not\textquotedblright{}, and \emph{apekṣā}, \textquotedblleft{}expectation\textquotedblright{}: \textquotedblleft{}not expected\textquotedblright{}, so \textquotedblleft{}unusual\textquotedblright{}.
\end{cousinweb}

\begin{grammarlens}[title={What you've built}]
One more word for this chapter.
\end{grammarlens}

\subsection*{Your turn}
\begin{itemize}
\raggedright
  \item \emph{Say it:} \emph{anokho}
  \item \emph{Say it:} \emph{anokho}, once more
  \item \emph{Say it:} say \emph{anokho}
  \item \emph{Recall:} say the Marwadi for special, then the Marwadi for plain, then say \emph{anokho} again
\end{itemize}

\begin{tcolorbox}[breakable,colback=teal!4,colframe=teal!35!black,title={Before you move on}]
What does \mw{अनोखो} mean? (\textquotedblleft{}Unusual\textquotedblright{}.) Say it once more.
\end{tcolorbox}
